\centerline{\bf \Huge Входная олимпиада}


\begin{flushright}
        \scriptsize{Оставь телефон всяк сюда входящий}
\end{flushright}


\problem{1}  На доску выписали все пятизначные числа, в записи которых все цифры различны. С доски стирают наибольшее и наименьшее из написанных до тех пор, пока не останутся два числа. Какие числа останутся на доске?

\problem{2} Моторная лодка плыла по реке. В полдень капитан выронил за борт шляпу. В час дня он 
обнаружил это и повернул обратно. В каком часу лодка поравняется со шляпой?

\problem{3} Сколько клеток пересекает диагональ в клетчатом прямоугольнике $25 * 36$?

\problem{4} В примере на сложение, цифры заменили буквами (одинаковые-одинаковыми, а разные-разными). 
Вот что получилось:БУЛОК + БЫЛО = МНОГО. Сколько же было булок? Их количество есть максимально возможное число МНОГО. 

\problem{5} В группе школьников у каждого хотя бы три знакомых среди членов этой группы. Все знакомые Фёдора незнакомы между собой, а все знакомые Баты знакомы между собой. Знакомы ли Фёдор и Бата?

\problem{6} Из клетчатого квадрата размером 7×7 по границам клеток вырезали равное количество квадратов 2×2 и прямоугольников 1×4. Какое наибольшее количество этих фигурок могло быть вырезано?
